\section{Plugin Architecture}
\label{sec:PluginArchitecture}

Durga supports a plugin-driven data pipeline model where BPMN \texttt{ServiceTask}
elements can be annotated with \texttt{camunda:plugin} attributes that resolve to
registered, governed data processing components.

\subsection{Plugin registry}

Plugins are defined by YAML descriptors under \filepath{plugins/} and indexed by
\filepath{plugins/catalog.yml}. Each descriptor declares:

\begin{itemize}
\item \texttt{id} --- unique plugin identifier (used in \texttt{camunda:plugin}),
\item \texttt{name} / \texttt{version} --- human-readable identifier,
\item \texttt{category} --- governed taxonomy (transform, validate, enrich, route, aggregate, sink),
\item \texttt{status} --- lifecycle state (experimental, stable, deprecated, retired),
\item \texttt{implementation} --- Java class implementing
\filepath{org.gautelis.durga.plugins.Plugin},
\item \texttt{config} --- typed configuration schema (field types, defaults, allowed values).
\end{itemize}

\subsection{Governed categories}

Each plugin category enforces a contract: topic naming, lifecycle events, and error
channel conventions. The scaffolder validates that generated executors follow the
category contract.

\begin{longtable}{|>{\raggedright\arraybackslash}p{0.22\textwidth}|>{\raggedright\arraybackslash}p{0.70\textwidth}|}
\hline
\rowcolor[gray]{0.9}\textbf{Category} & \textbf{Contract} \\
\hline
\textbf{transform} & 1:1 input$\to$output, ACTIVITY\_COMPLETED, dlq on failure \\
\textbf{validate} & Output + invalid branch, ACTIVITY\_ESCALATED on invalid \\
\textbf{enrich} & Input$\to$enriched output, may block on external lookup \\
\textbf{route} & Input$\to$N downstream channels, GATEWAY\_TAKEN lifecycle \\
\textbf{aggregate} & Windowed counting, ACTIVITY\_ENTERED on window start \\
\textbf{sink} & Terminal output, no downstream routing \\
\hline
\end{longtable}

\subsection{Plugin interface}

All plugins implement \filepath{org.gautelis.durga.plugins.Plugin}:

\begin{quote}
\shellcmd{String execute(String payload, String config) throws Exception}
\end{quote}

Generated executor classes instantiate the plugin at field level and call
\texttt{plugin.execute(payload, config)} for each message. The config string
comes from the BPMN \texttt{camunda:pluginConfig} property.

\subsection{Reference implementations}

The repository ships with five reference plugins that form the initial
standard library:

\begin{description}
\item[JSON Transform] Dot-notation field remapping.
\item[Field Filter] Whitelist/blacklist with flatten.
\item[KV Enricher] Inline key-value lookup and merge.
\item[Field Router] Value-based routing to named channels.
\item[Window Counter] Time-window counting with grouping.
\end{description}

\subsection{BPMN binding}

Plugin-annotated tasks use Camunda extension properties on \texttt{ServiceTask}
elements:

\begin{verbatim}
<bpmn:serviceTask id="transform" name="Transform Data">
  <bpmn:extensionElements>
    <camunda:properties>
      <camunda:property name="plugin" value="json-transform" />
      <camunda:property name="pluginConfig" value="name, email" />
    </camunda:properties>
  </bpmn:extensionElements>
</bpmn:serviceTask>
\end{verbatim}

Without a \texttt{camunda:plugin} attribute, \texttt{ServiceTask} elements fall
back to the existing generic worker generation.

\subsection{Generated executor}

The scaffolder resolves the plugin descriptor, validates it against the registry,
and generates a concrete executor class using the \texttt{pluginExecutorClass}
template. The executor:

\begin{itemize}
\item imports the plugin implementation class and the \texttt{Plugin} interface,
\item instantiates the plugin as a private field,
\item delegates each incoming message to \texttt{plugin.execute()},
\item emits \texttt{ACTIVITY\_COMPLETED} on success,
\item routes failures to a dead-letter topic (\texttt{\%task\%\_dlq}) with structured error records,
\item respects \texttt{ScopeCancellationRegistry} for interrupting boundary support.
\end{itemize}

\subsection{Current limitations}

\begin{itemize}
\item Plugin descriptors are read from the filesystem or classpath at scaffold time --- there is no runtime plugin discovery or hot-reload.
\item Config validation at scaffold time is not yet implemented --- invalid configs are caught at runtime by the generated executor.
\item The plugin catalog must be maintained manually; there is no registry publishing or version-resolution tooling.
\end{itemize}
